\chapter{绪\quad 论}

\section{引言}
然而混合动力驱动车辆要获得与单一热机驱动车辆相近的牵引特性，就要供给驱动车轮相近的能量。在混合动力驱动系统中，由于降低了主能源（由热机燃烧获得）的输出功率，为达到能量平衡，必须采用辅助能源补偿能量差值。目前采用的是相互独立的牵引蓄电池、高速飞轮以及超级电容等作为辅助能源，由于这种体系存在一定的缺陷，如控制复杂、比能量小等，使得其应用效果不明显，基于这一现状提出了本课题。
\section{研发背景及意义}
随着人类近代文明的进步，汽车已成为人类日常生活一种不可缺少的运输与代步工具，但它的大量存在也给人类生存带来了日益严峻的能源与环境问题。面对这一关系到人类生存环境的挑战，世界各国都在大投入着手开发一种低排放、高能效的车辆动力系统，从目前的研发进程来看主要呈现两种趋势：…… 
\section{单体高性能辅助发动机基本设计思路及关键技术}
\subsection{单体高性能辅助发动机的功能}
表1.1给出了不同……